\documentclass[10pt,a4paper]{amsart}
\usepackage{amsmath,amssymb,amsthm,graphicx}
\usepackage[top=1in,bottom=1in,right=1in,left=1in]{geometry}
\usepackage{amssymb,amsmath,amsthm}
\usepackage{pgfplots}
\usepackage{pgfplotstable}
\usepackage{arydshln}
\usepackage{verbatim}
\usepackage{titlesec}
\usepackage{subfigure}
\usepackage[linesnumbered,ruled]{algorithm2e}
\usepackage{bm}
\usepackage{pdfpages}
\usepackage{enumerate} 
\setlength{\parskip}{0.2\baselineskip}

\definecolor{forestgreen}{rgb}{0.13, 0.55, 0.13}
\usepackage[colorlinks, linkcolor=blue,citecolor=forestgreen,urlcolor = black, bookmarks=true]{hyperref}
\usepackage[T1]{fontenc}
\newcommand{\Mod}[1]{\ (\mathrm{mod}\ #1)}
 
\usepackage[nameinlink, noabbrev, capitalize]{cleveref}
\usepackage{tikz}
%\usepackage{libertine}
%\usepackage{txfonts}
%\usepackage{newtxtext,newtxmath}
\usepackage{times}

\pgfplotsset{compat=1.18}
\newcommand\blfootnote[1]{%
  \begingroup
  \renewcommand\thefootnote{}\footnote{#1}%
  \addtocounter{footnote}{-1}%
  \endgroup
}

\linespread{1.1}



\DeclareMathOperator{\Tr}{Tr}


\crefname{equation}{}{}
\crefrangeformat{equation}{(#3#1#4) --~(#5#2#6)}
\crefmultiformat{equation}{(#2#1#3)}{, (#2#1#3)}{, (#2#1#3)}{, (#2#1#3)}
\crefname{lem}{Lemma}{Lemmas}
\crefname{section}{Section}{Sections}
\crefname{subsubsubsection}{Section}{Sections}
\crefname{rem}{Remark}{Remarks}
\crefname{figure}{Figure}{Figures}
\crefname{table}{Table}{Tables}
\Crefname{lem}{Lemma}{Lemmas}
\crefname{thm}{Theorem}{Theorems}
\Crefname{thm}{Theorem}{Theorems}

\usepackage{commath}



\newtheorem{theorem}{Theorem}[section]
\newtheorem{lemma}{Lemma}[section]
\newtheorem{proposition}{Proposition}[section]
\newtheorem{corollary}{Corollary}[section]
\newtheorem{assumption}{Assumption}
\newtheorem{fact}{Fact}[section]
\newtheorem{claim}{Claim}
\newtheorem{question}{Question}
\newtheorem{remark}{Remark}
\newtheorem{definition}{Definition}[section]
\crefname{originallemma}{original lemma}{original lemmas}
\Crefname{originallemma}{Original lemma}{Original lemmas}
\makeatletter
\newlabel{lem:gibbs_comp}{{2.1}{source}{}{sourcepointer-1}{}}
\newlabel{lem:gibbs_comp@cref}{{[originallemma][0][]2.1}{[1][0][]source}{}{}{}}
\makeatother
\title{Contraction of interleaved projection scans from one contracting ordering}
\author{Jason Gaitonde, Elchanan Mossel}
\date{arxiv:2410.11136v2; CC BY 4.0}
\usepackage[UTF8,fontset=fandol]{ctex}
\begin{document}
\maketitle
\noindent\textit{Editorial note.} 原文对Lemma 2.1的引用只说明记号相似；本题所需Q与损失序列均在所附证明中重新定义并推出，无需把该引理结论当作先修。完整原论文随来源包提供。\par\medskip
\noindent\textit{Editorial note.} 原文第777行在文字明确调用Cauchy-Schwarz后仍使用等号连接到平方根上界；作者公式原样保留，此处作为排印疑点明确提示，独立审阅须裁定该文字与后继不等式是否已足以明确本步。未通过独审前不准入。\par\medskip
\noindent\textit{Editorial note.} 按定理前作者原段落，delta表示嵌入排列的谱隙，按非负谱隙参数理解；本题不将delta为任意负实数纳入结论。该作者语境一并保留。\par\medskip
\noindent\textit{Allowed standard prerequisites (editorial).}\begin{itemize}
\item Hilbert spaces, bounded operators and orthogonal projections
\item Pythagorean identity, triangle and Cauchy-Schwarz inequalities
\item Operator norm and elementary scalar inequalities
\item delta is the nonnegative spectral-gap parameter referred to in the author paragraph at source line722, not an arbitrary real number.
\end{itemize}\medskip
\ifdefined\hypertarget\hypertarget{sourcepointer-1}{}\fi
\noindent\textit{Source reference.} Original lemma 2.1: sources/projection-scans-arxiv2410.11136v2/source0.tex, source lines 396--433. Notation-only reference. The counted proof itself defines Q and its squared-norm losses and proves the Pythagorean telescoping identity; it does not invoke the lemma conclusion.\par\medskip
We now turn to proving our main results via abstract linear-algebraic bounds on operator norms. Throughout this section, for a bounded linear operator $\mathsf{A}:\mathcal{H}\to \mathcal{H}$ on a Hilbert space $\mathcal{H}$, we will write
\begin{equation*}
    \|\mathsf{A}\|_{\mathsf{op}} =\sup_{f\in \mathcal{H}} \frac{\|\mathsf{A}f\|}{\|f\|},
\end{equation*}
where $\|\cdot\|$ is the norm on $\mathcal{H}$ induced by the inner product.

\begin{theorem}
\label{thm:gen_gap}
        Let $\mathcal{H}$ be a Hilbert space and suppose $\mathsf{P}_1,\ldots, \mathsf{P}_n$ are orthogonal projections with respect to the inner product of $\mathcal{H}$ satisfying
        \begin{equation}
    \left\|\frac{1}{n}\sum_{i=1}^n \mathsf{P}_i\right\|_{\mathsf{op}}\leq 1-\delta.
\end{equation} 
Let $\mathsf{P}_{i_L}\ldots \mathsf{P}_{i_1}$ be any product of the orthogonal projection operators satisfying the following condition. For each $j\in [n]$, suppose there exists a minimal index $k_j\leq L$ satisfying $i_{k_j}=j$. We further assume these indices satisfy\footnote{In words, the indices $i_1,\ldots,i_L$ contain $[n]$ in ascending order of first appearance and are a minimal subsequence with this property.}:
\begin{equation}
\label{eq:order}
        1=k_1\leq k_2\leq\ldots\leq k_n=L.
\end{equation}
Then, it holds that
\begin{equation*}
    \|\mathsf{P}_{i_L}\ldots \mathsf{P}_{i_1}\|_{\mathsf{op}}^2\leq 1-\frac{n\delta}{8\sum_{j=1}^n k_j}.
\end{equation*}
\end{theorem}

\begin{proof}[Proof of \Cref{thm:gen_gap}]
We use similar notation as in \Cref{lem:gibbs_comp}. Fix any $f\in \mathcal{H}$ satisfying $\|f\|=1$ and define the sequence $0:=\gamma_0\leq \gamma_1\leq\ldots\leq \gamma_L$ via
\begin{equation*}
    \|\mathsf{P}_{i_{\ell}}\ldots \mathsf{P}_{i_1}f\|^2=1-\gamma_{\ell}.
\end{equation*}
For convenience of notation, further define the operators $\mathsf{Q}_{\ell}$ via 
\begin{equation*}
    \mathsf{Q}_{\ell} = \mathsf{P}_{i_{\ell}}\ldots \mathsf{P}_{i_1}.
\end{equation*}

For each $\ell\leq L-1$, the orthogonal decomposition
    \begin{equation*}
        \mathsf{Q}_{{\ell}}f = \mathsf{P}_{i_{\ell+1}}( \mathsf{Q}_{{\ell}}f)+(\mathsf{Q}_{{\ell}}f-P_{i_{\ell+1}}\mathsf{Q}_{{\ell}} f)= \mathsf{Q}_{{\ell+1}}f+(\mathsf{Q}_{{\ell}}f-\mathsf{Q}_{{\ell+1}} f)
    \end{equation*}
    certifies that
    \begin{equation}
    \label{eq:telescope}
        \|\mathsf{Q}_{{\ell}}f-\mathsf{Q}_{{\ell+1}} f\|^2=\|\mathsf{Q}_{{\ell}}f\|^2-\|\mathsf{Q}_{{\ell+1}}f\|^2=\gamma_{\ell+1}-\gamma_{\ell}.
    \end{equation}

    We now turn to bounding $\|\mathsf{P}_j f-f\|$. Since $i_{k_j}=j$ by definition,
    \begin{align*}
        \mathsf{P}_j f-f &= \mathsf{P}_{j}(f-\mathsf{P}_{i_{k_{j}-1}}\ldots \mathsf{P}_{i_1} f) + (\mathsf{P}_{i_{k_j}}\ldots \mathsf{P}_{i_1}f -f)\\
        &= \mathsf{P}_{j}(f-\mathsf{Q}_{{k_{j}-1}} f) + (\mathsf{Q}_{{k_j}}f -f)\\
        &=\mathsf{P}_{j}\left(\sum_{\ell=1}^{k_j-1} (\mathsf{Q}_{{\ell-1}} f-\mathsf{Q}_{{\ell}} f)\right)+\sum_{\ell=1}^{k_j} (\mathsf{Q}_{{\ell}} f -  \mathsf{Q}_{{\ell-1}} f).
    \end{align*}

    Taking norms, applying the triangle inequality and using the fact $\mathsf{P}_j$ is a projection, we obtain
\begin{align*}
    \| \mathsf{P}_j f-f\|&\leq \left\| \mathsf{P}_{j}\left(\sum_{\ell=1}^{k_j-1} (\mathsf{Q}_{{\ell-1}} f-\mathsf{Q}_{{\ell}} f)\right)\right\|+\sum_{\ell=1}^{k_j} \left\|\mathsf{Q}_{{\ell}} f -  \mathsf{Q}_{{\ell-1}} f\right\|\\
    &\leq  \sum_{\ell=1}^{k_j-1} \left\|\mathsf{Q}_{{\ell-1}} f-\mathsf{Q}_{{\ell}} f\right\|+\sum_{\ell=1}^{k_j} \left\|\mathsf{Q}_{{\ell}} f -  \mathsf{Q}_{{\ell-1}} f\right\|\\
    &\leq 2 \sum_{\ell=1}^{k_j} \left\|\mathsf{Q}_{{\ell}} f -  \mathsf{Q}_{{\ell-1}} f\right\|.
\end{align*}
Taking squares, applying Cauchy-Schwarz, and telescoping via \Cref{eq:telescope}, we find that
\begin{align*}
    \| \mathsf{P}_j f-f\|^2&\leq \left(2 \sum_{\ell=1}^{k_j} \left\|\mathsf{Q}_{{\ell}} f -  \mathsf{Q}_{{\ell-1}} f\right\|\right)^2\\
    &\leq 4\left(\sqrt{k_j}\sqrt{\sum_{\ell=1}^{k_j} \left\|\mathsf{Q}_{{\ell}} f -  \mathsf{Q}_{{\ell-1}} f\right\|^2}\right)^2\\
    &=4k_j\sum_{\ell=1}^{k_j} \left\|\mathsf{Q}_{{\ell}} f -  \mathsf{Q}_{{\ell-1}} f\right\|^2\\
    &=4k_j\sum_{\ell=1}^{k_j} (\gamma_{\ell}-\gamma_{\ell-1})\\
    &=4k_j\gamma_{k_j}.
\end{align*}

Finally, we observe by properties of orthogonal projection and the assumption $\|f\|=1$ that
\begin{align*}
    \left\| \frac{1}{n}\sum_{j=1}^n \mathsf{P}_jf\right\|^2&=\left\|f+ \frac{1}{n}\sum_{j=1}^n (\mathsf{P}_jf-f)\right\|^2\\
    &=\|f\|^2+2\left\langle f,\frac{1}{n}\sum_{j=1}^n (\mathsf{P}_jf-f)\right\rangle+\left\|\frac{1}{n}\sum_{j=1}^n (\mathsf{P}_jf-f)\right\|^2\\
    &\geq \|f\|^2-\frac{2}{n}\sum_{j=1}^n \vert \langle f,\mathsf{P}_jf-f\rangle\vert\\
    &=1-\frac{2}{n}\sum_{j=1}^n \|\mathsf{P}_jf-f\|^2\\
    &\geq 1-\frac{8}{n}\sum_{j=1}^n k_j\gamma_{k_j}\\
    &\geq 1-\frac{8\gamma_{L}}{n}\cdot \sum_{j=1}^n k_j.
\end{align*}
To see the last equality, note that
\begin{equation*}
    \langle f,\mathsf{P}_jf-f\rangle = \langle f-\mathsf{P}_jf,\mathsf{P}_jf-f\rangle+ \langle\mathsf{P}_jf,\mathsf{P}_jf-f\rangle = -\|\mathsf{P}_jf-f\|^2,
\end{equation*}
since any orthogonal projection is orthogonal to the remainder by the Pythagorean property.
Finally, by assumption the left hand side is at most $(1-\delta)^2\leq 1-\delta$ so rearranging yields
\begin{equation*}
    \gamma_{L} \geq \frac{n\delta}{8\sum_{j=1}^n k_j}.
\end{equation*}
Since this holds for any $f\in \mathcal{H}$ of unit norm, we have proven the claim.
\end{proof}

We will derive these results as a consequence of a general linear-algebraic result relating products of projections to that of an embedded permutation of the projections. Our assumption is that this embedded permutation has a spectral gap, and therefore there is an explicit spectral gap of the larger product that depends only on the original gap and the number of products in the sequence. The argument is conceptually straightforward: if the larger sequence of projections does not behave similarly to the single embedded permutation, this must be because the extra applications of projections noticeably moved the vector. But since any projection that moves a vector must significantly decrease the norm, this movement can also be charged to the spectral gap, and therefore one can obtain a polynomial relationship between the spectral gaps.

\begin{theorem}
\label{thm:gen_gap_2}
        Let $\mathcal{H}$ be a Hilbert space and suppose $\mathsf{P}_1,\ldots, \mathsf{P}_n$ are orthogonal projections with respect to the inner product of $\mathcal{H}$. Suppose that
        \begin{equation}
    \left\|\mathsf{P}_n\ldots \mathsf{P}_1\right\|_{\mathsf{op}}\leq 1-\delta.
\end{equation} 
Let $\mathsf{P}_{i_L}\ldots \mathsf{P}_{i_1}$ be any product of these orthogonal projection operators satisfying the following condition. For each $j\in [n]$, suppose there exists a minimal index $k_j\leq L$ satisfying $i_{k_j}=j$. We further assume these indices satisfy:
\begin{equation}
\label{eq:order_2}
        1=k_1\leq k_2\leq\ldots\leq k_n=L.
\end{equation}
Then, it holds that
\begin{equation*}
    \|\mathsf{P}_{i_L}\ldots \mathsf{P}_{i_1}\|_{\mathsf{op}}\leq 1-\frac{\delta^2}{8(L-n+1)}.
\end{equation*}
\end{theorem}

\begin{proof}
    We will provide two lower bounds on the spectral gap and show that at least one gives the desired value. As in \Cref{thm:gen_gap}, fix any $f\in \mathcal{H}$ satisfying $\|f\|=1$, set $\mathsf{Q}_{\ell} = P_{i_{\ell}}\cdots P_{i_1}$ and define the sequence $0=\lambda_0\leq \ldots\leq \lambda_L$ via
\begin{equation*}
    \|\mathsf{Q}_{{\ell}}f\|^2 = 1-\lambda_{\ell}.
\end{equation*}
Our goal is to lower bound $\lambda_L$. An identical argument to \Cref{thm:gen_gap} by writing out the orthogonal decompositions implies that for any $\ell$,
\begin{equation}
\label{eq:telescope_2}
    \|\mathsf{Q}_{{\ell+1}}f - \mathsf{Q}_{{\ell}}f\|^2 = \lambda_{\ell+1} - \lambda_{\ell}.
\end{equation}
On the one hand, for any subset $S$ of $[L]$, we trivially have by nonnegativity of the summands that
\begin{equation}
\label{eq:bound_1}
    \lambda_L = \sum_{\ell=1}^L (\lambda_{\ell}-\lambda_{\ell-1})\geq \sum_{\ell\in S} (\lambda_{\ell}-\lambda_{\ell-1}).
\end{equation}

On the other hand, since $i_L=n$ by assumption, we have
\begin{align*}
    \mathsf{Q}_L f &=\mathsf{P}_{i_{L}}\ldots \mathsf{P}_{i_1}f\\
    &= \mathsf{P}_n \mathsf{Q}_{{L-1}}f\\
    &=\mathsf{P}_n\mathsf{P}_{n-1}\ldots \mathsf{P}_1f + \mathsf{P}_n(\mathsf{Q}_{{L-1}}f-\mathsf{P}_{n-1}\ldots \mathsf{P}_1f).
\end{align*}
Taking norms, the triangle inequality, and the fact projections cannot increase norm imply that
\begin{align}
    \|\mathsf{Q}_{{L}}f\|&\leq \|\mathsf{P}_n\mathsf{P}_{n-1}\ldots \mathsf{P}_1f\| + \|\mathsf{P}_n(\mathsf{Q}_{{L-1}}f-\mathsf{P}_{n-1}\ldots \mathsf{P}_1f)\| \nonumber\\
    &\leq 1-\delta + \|\mathsf{Q}_{{L-1}}f-\mathsf{P}_{n-1}\ldots \mathsf{P}_1f\| \label{eq:initial}.
\end{align}
We now treat the last term similarly as before: write
\begin{equation*}
    \mathsf{Q}_{{L-1}}f = \mathsf{Q}_{{k_{n-1}}}f+\sum_{\ell=k_{n-1}}^{L-2} (\mathsf{Q}_{{\ell+1}}f-\mathsf{Q}_{{\ell}}f) ,
\end{equation*}
so a similar set of manipulations with the triangle inequality, Cauchy-Schwarz, and the fact projections can only decrease distance (noting $i_{k_{n-1}}=n-1$ by definition), and \Cref{eq:telescope_2} yields
\begin{align*}
    \|\mathsf{Q}_{{L-1}}f-\mathsf{P}_{n-1}\ldots \mathsf{P}_1f\|&=\left\| \mathsf{Q}_{{k_{n-1}}}f-\mathsf{P}_{n-1}\ldots \mathsf{P}_1f+\sum_{\ell=k_{n-1}}^{L-2} (\mathsf{Q}_{{\ell+1}}f-\mathsf{Q}_{{\ell}}f)\right\|\\
    &\leq \| \mathsf{Q}_{{k_{n-1}}}f-\mathsf{P}_{n-1}\ldots \mathsf{P}_1f\|+\sum_{\ell=k_{n-1}}^{L-2} \|\mathsf{Q}_{{\ell+1}}f-\mathsf{Q}_{{\ell}}f\|\\
    &=\| \mathsf{P}_{n-1}\mathsf{Q}_{{k_{n-1}-1}}f-\mathsf{P}_{n-1}\ldots \mathsf{P}_1f\|+\sqrt{k_{n}-k_{n-1}-1} \sqrt{\sum_{\ell=k_{n-1}}^{L-2} \|\mathsf{Q}_{{\ell+1}}f-\mathsf{Q}_{{\ell}}f\|^2}\\
    &\leq \|\mathsf{Q}_{{k_{n-1}-1}}f-\mathsf{P}_{n-2}\ldots \mathsf{P}_1f\|+\sqrt{k_n-k_{n-1}-1}\sqrt{\sum_{\ell=k_{n-1}}^{L-2} (\lambda_{\ell+1}-\lambda_{\ell})}.
\end{align*}

Proceeding in the same way recursively to deal with the first term (and using $i_1=1$ at the end of the recursion), we then obtain by Cauchy-Schwarz that
\begin{align*}
    \|\mathsf{Q}_{{L-1}}f-\mathsf{P}_{n-1}\ldots \mathsf{P}_1f\|&\leq \sum_{j=1}^{n-1} \sqrt{k_{j+1}-k_{j}-1}\sqrt{\sum_{\ell=k_{j}}^{k_{j+1}-2} (\lambda_{\ell+1}-\lambda_{\ell})} \nonumber\\
    &\leq \sqrt{\sum_{j=1}^{n-1} (k_{j+1}-k_j-1)}\sqrt{\sum_{\ell\in S}(\lambda_{\ell+1}-\lambda_{\ell})}\nonumber\\
    &=\sqrt{L-n}\sqrt{\sum_{\ell\in S}(\lambda_{\ell+1}-\lambda_{\ell})},
\end{align*}
where $S$ is a suitable subset of $[L]$; note that we use the fact that the inner summands for the increments are disjoint across different $k$ by construction. Putting this back with \Cref{eq:initial}, we have
\begin{equation}
\label{eq:bound_2}
    \|\mathsf{P}_{i_{L}}\ldots \mathsf{P}_{i_1}f\|\leq 1-\delta+ \sqrt{L-n}\sqrt{\sum_{\ell\in S}(\lambda_{\ell+1}-\lambda_{\ell})}
\end{equation}

For this subset $S$, let 
\begin{equation*}
    \kappa = \sum_{\ell\in S}(\lambda_{\ell+1}-\lambda_{\ell}).
\end{equation*}
By \Cref{eq:bound_1} and \Cref{eq:bound_2}, we observe that
\begin{equation*}
    \lambda_L \geq \max\left\{\kappa,\delta- \sqrt{L-n}\kappa^{1/2} \right\}.
\end{equation*}
It is not difficult to see that the right hand side  has value at least $\delta^2/4(L-n+1)$ by a simple case analysis (note  though that the bound improves to $\delta$ when $L=n$, as it should). Therefore, we conclude that 
\begin{equation*}
    \lambda_L\geq \frac{\delta^2}{4(L-n+1)}
\end{equation*}
as claimed, which yields the desired operator norm bound (without the square) by Bernoulli's inequality.
\end{proof}

\end{document}
